% Options for packages loaded elsewhere
\PassOptionsToPackage{unicode}{hyperref}
\PassOptionsToPackage{hyphens}{url}
\documentclass[
  12pt,
]{ctexart}
\usepackage{xcolor}
\usepackage{amsmath,amssymb}
\setcounter{secnumdepth}{-\maxdimen} % remove section numbering
\usepackage{iftex}
\ifPDFTeX
  \usepackage[T1]{fontenc}
  \usepackage[utf8]{inputenc}
  \usepackage{textcomp} % provide euro and other symbols
\else % if luatex or xetex
  \usepackage{unicode-math} % this also loads fontspec
  \defaultfontfeatures{Scale=MatchLowercase}
  \defaultfontfeatures[\rmfamily]{Ligatures=TeX,Scale=1}
\fi
\usepackage{lmodern}
\ifPDFTeX\else
  % xetex/luatex font selection
\fi
% Use upquote if available, for straight quotes in verbatim environments
\IfFileExists{upquote.sty}{\usepackage{upquote}}{}
\IfFileExists{microtype.sty}{% use microtype if available
  \usepackage[]{microtype}
  \UseMicrotypeSet[protrusion]{basicmath} % disable protrusion for tt fonts
}{}
\makeatletter
\@ifundefined{KOMAClassName}{% if non-KOMA class
  \IfFileExists{parskip.sty}{%
    \usepackage{parskip}
  }{% else
    \setlength{\parindent}{0pt}
    \setlength{\parskip}{6pt plus 2pt minus 1pt}}
}{% if KOMA class
  \KOMAoptions{parskip=half}}
\makeatother
\usepackage{longtable,booktabs,array}
\usepackage{calc} % for calculating minipage widths
% Correct order of tables after \paragraph or \subparagraph
\usepackage{etoolbox}
\makeatletter
\patchcmd\longtable{\par}{\if@noskipsec\mbox{}\fi\par}{}{}
\makeatother
% Allow footnotes in longtable head/foot
\IfFileExists{footnotehyper.sty}{\usepackage{footnotehyper}}{\usepackage{footnote}}
\makesavenoteenv{longtable}
\setlength{\emergencystretch}{3em} % prevent overfull lines
\providecommand{\tightlist}{%
  \setlength{\itemsep}{0pt}\setlength{\parskip}{0pt}}
\usepackage{bookmark}
\IfFileExists{xurl.sty}{\usepackage{xurl}}{} % add URL line breaks if available
\urlstyle{same}
\hypersetup{
  hidelinks,
  pdfcreator={LaTeX via pandoc}}

\author{SCX}
\date{2026}

\begin{document}

\section{SCX Self-Evolution Theory ---
Mathematical Formalization}\label{scx-ux81eaux6211ux8fdbux5316ux7406ux8bba-ux6570ux5b66ux5f62ux5f0fux5316}

\begin{quote}
Formalize the SCX self-evolution closed loop (judge -> store -> update -> re-judge -> re-update ->
...) as a rigorous mathematical theory, and establish connections with known theoretical frameworks.
\end{quote}

\textbf{Last Updated}: 2026-06-28

\begin{center}\rule{0.5\linewidth}{0.5pt}\end{center}

\subsection{Problem}\label{ux95eeux9898}

The SCX system forms the following closed loop during operation:

\begin{verbatim}
Judge (S_t scoring)  scoring) → Store (M_t  memory bank) → Update (S_{t+1}  update scoring) → Re-judge → ...
                                ↘ Update (θ_{t+1} NEP  student 更新) ↗
\end{verbatim}

\textbf{Core Problem}: Does this closed loop converge? To what? At what rate? Under what conditions does it diverge?

\textbf{Difficulty}：Gatekeeper \(S_t\) 和 NEP  student  \(f_{\\theta_t}\)
 form a coupled dynamical system------\(S_t\)  determines  \(f_{\\theta_t}\)
 the training distribution, and  \(f_{\\theta_t}\)  output provides feedback to  \(S_t\)。

\begin{center}\rule{0.5\linewidth}{0.5pt}\end{center}

\subsection{Core Insight}\label{ux6838ux5fc3ux6d1eux5bdf}

\begin{verbatim}
The correct mathematical object for SCX self-evolution is the coupled dynamical system (S_t, theta_t, M_t),
whose global behavior is governed by the monotonic descent property of the Lyapunov function V(S, theta).
\end{verbatim}

其中 \(S_t: \mathcal{X} \to [0,1]\) 是 gatekeeper 评分函数，\(\theta_t\)
是 NEP  student 参数，\(M_t\) 是单调增长的经验记忆库。

\begin{center}\rule{0.5\linewidth}{0.5pt}\end{center}

\subsection{Nine Core Files}\label{ux4e5dux4e2aux6838ux5fc3ux6587ux4ef6}

\begin{longtable}[]{@{}
  >{\raggedright\arraybackslash}p{(\linewidth - 6\tabcolsep) * \real{0.2143}}
  >{\raggedright\arraybackslash}p{(\linewidth - 6\tabcolsep) * \real{0.2143}}
  >{\raggedright\arraybackslash}p{(\linewidth - 6\tabcolsep) * \real{0.3571}}
  >{\raggedright\arraybackslash}p{(\linewidth - 6\tabcolsep) * \real{0.2143}}@{}}
\toprule\noalign{}
\begin{minipage}[b]{\linewidth}\raggedright
File
\end{minipage} & \begin{minipage}[b]{\linewidth}\raggedright
Content
\end{minipage} & \begin{minipage}[b]{\linewidth}\raggedright
Core Result
\end{minipage} & \begin{minipage}[b]{\linewidth}\raggedright
Status
\end{minipage} \\
\midrule\noalign{}
\endhead
\bottomrule\noalign{}
\endlastfoot
\texttt{01\_symbol\_system.md} & Symbol System and Problem Setup & 结构空间
\(\mathbb{X}\)、Status空间 \(\mathbb{Z}\)、更新算子 \(\Phi\)、新Assumptions SE-A1
 to  SE-A6 & \textbf{Complete} \\
\texttt{02\_dynamical\_system.md} & Discrete Dynamical System Formalization & Lyapunov 函数
\(V(z_t)\)、不动点存在性(Theorem  4-5)、单调下降(Theorem  6)、相图四种区域 &
\textbf{Complete} \\
\texttt{03\_online\_learning\_regret.md} & Online Learning Regret  Analysis & Regret
\(\leq 2GR\sqrt{T}\) (OGD)、 delayed feedback Regret
\(\leq 2GR\sqrt{T} + G^2 D_{\max}\sqrt{T}\) (Theorem 8) & \textbf{Complete} \\
\texttt{04\_bayesian\_update.md} & Bayesian Update Interpretation & \(S_t\)
is 后验均值、贝叶斯鞅 \(\mathbb{E}[S_{t+1}|\mathcal{F}_t]=S_t\)、Doob
收敛 \(S_t \to S_\infty\) a.s. & \textbf{Complete} \\
\texttt{05\_stochastic\_approximation.md} & Stochastic Approximation Analysis & Robbins-Monro
形式 \(\theta_{t+1} = \theta_t - \alpha_t \nabla \ell + \xi_t\)、ODE
方法、双时间尺度 Analysis & \textbf{Complete} \\
\texttt{06\_fixed\_point\_convergence.md} &
\textbf{中心Theorem：不动点与收敛} & \textbf{Theorem SE-1}：has 限结构空间 +
Lipschitz + 充分退火 → \((S_t, \theta_t) \to (S^*, \theta^*)\)
a.s.；Four Convergence Paths & \textbf{Complete} \\
\texttt{07\_completeness.md} & Completeness Analysis & \textbf{Theorem
SE-2}：物理约束下存在has 限 \(T^*\) 使得 \(t \geq T^*\) 后系统达到
\(\varepsilon\)-近似不动点；Gödel  analogy & \textbf{Complete} \\
\texttt{08\_theory\_connections.md} & Connections to Known Theories & AlphaZero
self-play、贝叶斯优化、主动学习、Solomonoff 归纳的形式化对比表 &
\textbf{Complete} \\
\texttt{09\_verification\_report.md} & Verification Report &
一致性检查、跨Theorem交叉引用、Symbol/Concept冲突、Proof缺口（10 个)、open problems（5
个) & \textbf{Complete} \\
\end{longtable}

\begin{center}\rule{0.5\linewidth}{0.5pt}\end{center}

\subsection{Core Theorems}\label{ux6838ux5fc3ux5b9aux7406}

\subsubsection{Theorem SE-1: Convergence of SCX
Self-Evolution}\label{theorem-se-1-convergence-of-scx-self-evolution}

\textbf{Conditions (C1-C7)}： 1. 结构空间 \(\mathbb{X}\)
has 限（has 限 data、has 限精度) 2. NEP  student  \(f_\theta\) 关于 \(\theta\) 是
\(L_f\)-Lipschitz 的 3. Gatekeeper \(S_t\) 关于其参数是
\(L_S\)-Lipschitz 的 4. 学习率满足 Robbins-Monro
Conditions：\(\sum \alpha_t = \infty, \sum \alpha_t^2 < \infty\) 5. 给定
\(S_t\) 的Conditions下，记忆库采样是 i.i.d. 的 6.
充分退火：\(\alpha_t \to 0\)，且 gatekeeper  update frequency \(\to 0\) 7.
Gatekeeper
更新步长has 界：\(\|S_{t+1} - S_t\|_\infty \leq \eta_t\)，\(\sum \eta_t < \infty\)

\textbf{Conclusion}：序列 \((S_t, \theta_t)\) 几乎必然收敛到联合不动点
\((S^*, \theta^*)\)，其中： - \(S^*\)
自洽：\(S^*(x) = \mathbb{P}(\text{correct} \mid x, \mathcal{M}_\infty)\)
- \(\theta^*\) 是 \(S^*\) 诱导分布下 NEP 期望损失的局部极小值

\textbf{ProofPath}：Lemma SE-1.1 (Lyapunov 超鞅) + Lemma SE-1.2
(has 限类型) + Lemma SE-1.3 (步长消失) + Lemma SE-1.4 (极限点即不动点)。

\subsubsection{Theorem SE-2: Completeness
Bound}\label{theorem-se-2-completeness-bound}

在物理约束（has 限 data \(N_{\max}\)、has 限精度
\(\varepsilon_{\text{mach}}\)、has 限计算)下，存在has 限时间 \(T^*\)
使得对所has  \(t \geq T^*\)，系统处于 \(\varepsilon\)-近似不动点。

\begin{center}\rule{0.5\linewidth}{0.5pt}\end{center}

\subsection{Four Convergence Paths}\label{ux56dbux79cdux6536ux655bux8defux5f84}

\begin{longtable}[]{@{}
  >{\raggedright\arraybackslash}p{(\linewidth - 4\tabcolsep) * \real{0.3333}}
  >{\raggedright\arraybackslash}p{(\linewidth - 4\tabcolsep) * \real{0.3333}}
  >{\raggedright\arraybackslash}p{(\linewidth - 4\tabcolsep) * \real{0.3333}}@{}}
\toprule\noalign{}
\begin{minipage}[b]{\linewidth}\raggedright
Path
\end{minipage} & \begin{minipage}[b]{\linewidth}\raggedright
Characteristics
\end{minipage} & \begin{minipage}[b]{\linewidth}\raggedright
Conditions
\end{minipage} \\
\midrule\noalign{}
\endhead
\bottomrule\noalign{}
\endlastfoot
\textbf{I. Classical Convergence} & \((S_t, \theta_t) \to (S^*, \theta^*)\)， monotonic improvement
& C1-C7  all satisfied, sufficient annealing \\
\textbf{II. Limit Cycle} &  system oscillates periodically among finite configurations &  insufficient annealing, excessive coupling \\
\textbf{III. Perpetual Discovery} &  new structures continuously discovered, \(\mathcal{M}_t\)  grows without bound &
 open physical world, unlimited exploration budget \\
\textbf{IV. Divergent Collapse} & \(S_t\)  degenerates, quality declines &  feedback loop breaks, NEP
 feedback noise excessive \\
\end{longtable}

Path II-IV The precise boundary conditions are\textbf{open problems}（ see 
\texttt{09\_verification\_report.md})。

\begin{center}\rule{0.5\linewidth}{0.5pt}\end{center}

\subsection{Relationship with Existing Theoretical Frameworks}\label{ux4e0eux73b0ux6709ux7406ux8bbaux4f53ux7cfbux7684ux5173ux7cfb}

\subsubsection{Subordination Relationships}\label{ux4eceux5c5eux5173ux7cfb}

\begin{verbatim}
Theorem 3 (不可识别性) ───  provides the necessity argument for self-evolution ──→  the self-evolution closed loop must exist
        │
        ▼
Theorem 1 (噪声检测) ──── is  S_t  provides initialization guarantees ──→ Theorem SE-1  initial conditions
        │
        ▼
Theorem SE-1 (自进化收敛) ───  central result ──→  closed loop converges asymptotically to a self-consistent point
        │
        ├──→ Theorem 2 ( Weak Feature Bound) ──  limits improvement speed
        ├──→ Theorem 4' ( Minimax Optimality) ── 提供 S^* 的渐近Optimal 性
        └──→ Theorem SE-2 ( Completeness Bound) ──  physical finiteness enforces finite-time convergence
\end{verbatim}

\subsubsection{Newly Introduced Assumptions (SE-A1  to 
SE-A6)}\label{ux65b0ux5f15ux5165ux5047ux8bbe-se-a1-ux81f3-se-a6}

\begin{longtable}[]{@{}lll@{}}
\toprule\noalign{}
Assumptions & Content & 对应现has Assumptions \\
\midrule\noalign{}
\endhead
\bottomrule\noalign{}
\endlastfoot
\textbf{SE-A1} & Lyapunov 下降 & \emph{ Newly introduced, core assumption} \\
\textbf{SE-A2} & has 限结构空间 &  Physical constraint (non-statistical assumption) \\
\textbf{SE-A3} & Lipschitz NEP &  Smoothness assumption \\
\textbf{SE-A4} & Conditions i.i.d. 采样 & A1-A2  extension \\
\textbf{SE-A5} & Robbins-Monro 学习率 &  Standard  SA Conditions \\
\textbf{SE-A6} & Gatekeeper 更新has 界 &  Annealing condition \\
\end{longtable}

\begin{center}\rule{0.5\linewidth}{0.5pt}\end{center}

\subsection{Summary of Connections to Known Theories}\label{ux4e0eux5df2ux77e5ux7406ux8bbaux7684ux8fdeux63a5ux6458ux8981}

\begin{longtable}[]{@{}
  >{\raggedright\arraybackslash}p{(\linewidth - 4\tabcolsep) * \real{0.2273}}
  >{\raggedright\arraybackslash}p{(\linewidth - 4\tabcolsep) * \real{0.2273}}
  >{\raggedright\arraybackslash}p{(\linewidth - 4\tabcolsep) * \real{0.5455}}@{}}
\toprule\noalign{}
\begin{minipage}[b]{\linewidth}\raggedright
Theoretical Framework
\end{minipage} & \begin{minipage}[b]{\linewidth}\raggedright
Core Mechanism
\end{minipage} & \begin{minipage}[b]{\linewidth}\raggedright
与 SCX Key Differences from SCX Self-Evolution
\end{minipage} \\
\midrule\noalign{}
\endhead
\bottomrule\noalign{}
\endlastfoot
\textbf{AlphaZero} & Policy iteration + MCTS & AlphaZero
 generates synthetic games; SCX  relies on real  NEP  data \\
\textbf{贝叶斯优化} & Acquisition + Surrogate & BO  optimizes static black box; SCX
 target distribution evolves \\
\textbf{主动学习} & Query strategy + Oracle & AL 的 oracle  immediate oracle feedback; SCX
的 NEP  delayed feedback \\
\textbf{Solomonoff 归纳} & Universal prior & Solomonoff  uncomputable; SCX
 restricted to finite hypothesis class \\
\end{longtable}

See  \texttt{08\_theory\_connections.md}  for the complete comparison table.

\begin{center}\rule{0.5\linewidth}{0.5pt}\end{center}

\subsection{Theoretical Status Assessment}\label{ux7406ux8bbaux72b6ux6001ux8bc4ux4f30}

\begin{longtable}[]{@{}
  >{\raggedright\arraybackslash}p{(\linewidth - 4\tabcolsep) * \real{0.3333}}
  >{\raggedright\arraybackslash}p{(\linewidth - 4\tabcolsep) * \real{0.3333}}
  >{\raggedright\arraybackslash}p{(\linewidth - 4\tabcolsep) * \real{0.3333}}@{}}
\toprule\noalign{}
\begin{minipage}[b]{\linewidth}\raggedright
Category
\end{minipage} & \begin{minipage}[b]{\linewidth}\raggedright
Proportion
\end{minipage} & \begin{minipage}[b]{\linewidth}\raggedright
Description
\end{minipage} \\
\midrule\noalign{}
\endhead
\bottomrule\noalign{}
\endlastfoot
\textbf{Rigorous Theory} & \textasciitilde40\% & Theorem SE-1
 finite-state version、 Bayesian martingale convergence、Robbins-Monro  standard convergence、Regret  upper bound \\
\textbf{Formalized Conjectures} & \textasciitilde20\% &
 continuous state-space extension、 two-timescale convergence、KL  contraction rate \\
\textbf{Reasonable Assumptions} & \textasciitilde40\% & Exact  Lyapunov
 function form、 limit cycle conditions、Gödel  analogy、 optimal annealing strategy \\
\end{longtable}

\begin{center}\rule{0.5\linewidth}{0.5pt}\end{center}

\subsection{open problems（Priority Ordering)}\label{ux5f00ux653eux95eeux9898ux4f18ux5148ux7ea7ux6392ux5e8f}

\begin{longtable}[]{@{}lll@{}}
\toprule\noalign{}
Priority & Problem & File \\
\midrule\noalign{}
\endhead
\bottomrule\noalign{}
\endlastfoot
\textbf{P0} & Explicit  Lyapunov 函数 \(\Phi\)  definition and descent proof &
\texttt{06}, \texttt{09} \\
\textbf{P1} & Precise bounding of the four convergence paths & \texttt{06} \\
\textbf{P1} &  coupled  \((S_t, \theta_t)\)  system tight convergence rate & \texttt{05},
\texttt{06} \\
\textbf{P2} & Optimal  gatekeeper  update frequency & \texttt{02}, \texttt{03} \\
\textbf{P2} & \(\varepsilon\)-Minimum memory bank size for convergence & \texttt{07} \\
\textbf{P3} & Non-asymptotic (finite-time) guarantees & \texttt{03}, \texttt{06} \\
\textbf{P4} & Robustness under distribution shift & \texttt{05} \\
\textbf{P5} & Gödel  analogy rigorization & \texttt{07} \\
\end{longtable}

\begin{center}\rule{0.5\linewidth}{0.5pt}\end{center}

\subsection{File Structure}\label{ux6587ux4ef6ux7ed3ux6784}

\begin{verbatim}
theory/self_evolution/
├── README.md                              ← This file
├── 01_symbol_system.md                    # Symbol System and Problem Setup
├── 02_dynamical_system.md                 # Discrete Dynamical System Formalization
├── 03_online_learning_regret.md           # Online Learning Regret  Analysis
├── 04_bayesian_update.md                  # Bayesian Update Interpretation
├── 05_stochastic_approximation.md         # Stochastic Approximation Analysis
├── 06_fixed_point_convergence.md          # Fixed Point and Convergence Theorem (Center)
├── 07_completeness.md                     # Completeness Analysis
├── 08_theory_connections.md               # Connections to Known Theories
└── 09_verification_report.md              # Verification Report
\end{verbatim}

\begin{center}\rule{0.5\linewidth}{0.5pt}\end{center}

\subsection{Cross-Domain Applicability}\label{ux8de8ux9886ux57dfux9002ux7528}

\begin{longtable}[]{@{}
  >{\raggedright\arraybackslash}p{(\linewidth - 6\tabcolsep) * \real{0.1053}}
  >{\raggedright\arraybackslash}p{(\linewidth - 6\tabcolsep) * \real{0.1930}}
  >{\raggedright\arraybackslash}p{(\linewidth - 6\tabcolsep) * \real{0.3333}}
  >{\raggedright\arraybackslash}p{(\linewidth - 6\tabcolsep) * \real{0.3684}}@{}}
\toprule\noalign{}
\begin{minipage}[b]{\linewidth}\raggedright
Domain
\end{minipage} & \begin{minipage}[b]{\linewidth}\raggedright
\(S_t\) Role
\end{minipage} & \begin{minipage}[b]{\linewidth}\raggedright
\(f_{\theta_t}\) Role
\end{minipage} & \begin{minipage}[b]{\linewidth}\raggedright
\(\mathcal{M}_t\) Role
\end{minipage} \\
\midrule\noalign{}
\endhead
\bottomrule\noalign{}
\endlastfoot
\textbf{MLIP} & Judge  DFT  label correctness & NEP  potential energy surface model & Accumulated  DFT
 computation trajectories \\
\textbf{Medical Imaging} &  judge diagnostic label quality &  auxiliary diagnostic model &
 accumulated expert review records \\
\textbf{LLM} &  judge annotation/ generation quality &  downstream task model &  accumulated human review logs \\
\textbf{Autonomous Driving} &  judge detection box/ annotation quality &  perception model &
 accumulated road-test validation data \\
\end{longtable}

\begin{center}\rule{0.5\linewidth}{0.5pt}\end{center}

\subsection{References (Newly Added)}\label{ux53c2ux8003ux6587ux732eux65b0ux589e}

\begin{enumerate}
\def\labelenumi{\arabic{enumi}.}
\tightlist
\item
  Robbins, H., \& Monro, S. (1951). A stochastic approximation method.
  \emph{Annals of Mathematical Statistics}, 22(3), 400-407.
\item
  Borkar, V. S. (2008). \emph{Stochastic Approximation: A Dynamical
  Systems Viewpoint}. Cambridge University Press.
\item
  Kushner, H. J., \& Yin, G. G. (2003). \emph{Stochastic Approximation
  and Recursive Algorithms and Applications} (2nd ed.). Springer.
\item
  Doob, J. L. (1953). \emph{Stochastic Processes}. Wiley.
\item
  Bernardo, J. M., \& Smith, A. F. M. (1994). \emph{Bayesian Theory}.
  Wiley.
\item
  van der Vaart, A. W. (1998). \emph{Asymptotic Statistics}. Cambridge
  University Press.
\item
  Silver, D., et al.~(2017). Mastering the game of Go without human
  knowledge. \emph{Nature}, 550, 354-359.
\item
  Shahriari, B., et al.~(2016). Taking the human out of the loop: A
  review of Bayesian optimization. \emph{Proceedings of the IEEE},
  104(1), 148-175.
\item
  Settles, B. (2009). Active learning literature survey.
  \emph{University of Wisconsin-Madison Technical Report}.
\item
  Solomonoff, R. J. (1964). A formal theory of inductive inference.
  \emph{Information and Control}, 7(1), 1-22, 224-254.
\item
  Li, M., \& Vitányi, P. (2008). \emph{An Introduction to Kolmogorov
  Complexity and Its Applications} (3rd ed.). Springer.
\item
  Zinkevich, M. (2003). Online convex programming and generalized
  infinitesimal gradient ascent. \emph{ICML}.
\item
  Cesa-Bianchi, N., \& Lugosi, G. (2006). \emph{Prediction, Learning,
  and Games}. Cambridge University Press.
\item
  SCX Theorem 1-3. \texttt{../theorems/}.
\item
  SCX Unified Theorem Document. \texttt{../THEOREMS\_UNIFIED.md}.
\end{enumerate}

\begin{center}\rule{0.5\linewidth}{0.5pt}\end{center}

\emph{End of README.md --- SCX Self-Evolution Theory}

\end{document}
